% ===================================================================
%  도표 제 8 호 — 거두기.
%
%  Traduction, faite en 2026, en coreen d'aujourd'hui, sur l'ido de
%  texto/io. Regles de la colonne : voir 10-dopyo-01.tex. Deux
%  scenes, deux numerotations : les cles portent c1 et c2. Le
%  fac-simile compose « 13) » sans sa parenthese ouvrante ; la
%  traduction ecrit la reference bien formee, comme le font deja les
%  colonnes hindi, marathe et telougoue.
%
%  LE FLEAU (40), ET C'EST LA QUATRIEME LANGUE A LE RENCONTRER. Les
%  trois precedentes ne l'avaient pas : l'egyptien ecrivait عصي
%  الدراس, « les batons de battage », parce que l'Egypte bat au نورج ;
%  le marathi मळणीच्या काठ्या parce que le Maharashtra fait pietiner
%  les bœufs ; le telougou నూర్పిడి కర్రలు pour la meme raison. Trois
%  familles de langues, un seul trou, la meme periphrase. Le coreen a
%  le mot ET l'outil : 도리깨, deux batons relies par une lanière,
%  exactement l'objet de la planche. Le trou etait dans la CHOSE et
%  non dans les langues — et il fallait une quatrieme langue, qui a
%  la chose, pour achever de le prouver.
%
%  C'EST LE DEUXIEME OBJET DE CE TABLEAU QUI TOMBE AINSI, apres le
%  marron chaud du tableau 7. Deux fois de suite, la colonne coreenne
%  se trouve du bon cote d'une coupure que quatre colonnes avaient
%  documentee de l'autre.
%
%  ET LES TROIS OISEAUX DU TABLEAU 8 TELOUGOU ONT AUSSI LEUR NOM ICI :
%  할미새 la bergeronnette (19), 종달새 l'alouette (41), 제비
%  l'hirondelle (56), et j'ajoute 방울새 le chardonneret (36) et
%  멧도요 la becasse (61). La colonne telougoue avait note ces trois-la
%  comme non resolus faute d'oiseaux qui passent au-dessus de
%  l'Andhra ; ils passent au-dessus de la Coree.
%
%  RESTENT DEUX EMPRUNTS DE PLANTES, ET LA MEME REGLE QUE PARTOUT :
%  디기탈리스 la digitale (3) et 카나리아 le serin (35) ne vivent pas
%  ici. On ecrit l'emprunt. Le bleuet, lui, tombe juste — 수레국화 est
%  le mot coreen, et la fleur pousse.
%
%  VINGT-QUATRE PAIRES. Le tableau alterne le recit et la description,
%  et le compte s'en ressent : la moisson enumere des outils, la
%  vendange attache des objets a des objets.
% ===================================================================

%%K t08-noto-1 noto
\VUnotes{143.2mm}{%
\textit{(*) 이 낱말이 또렷하지 않기 때문이다. 아마 거두기인가, 포도 거두기인가,
사과 거두기인가? Mondo XIV, 242쪽에서 제안한 «} \VUgras{mesono}
\textit{»를 쓰는 것이 나을지도 모른다. 그러나 지금까지 그 새 어근은 아카데미가
받아들이지 않았고, 이도의 관례에 따라 논의를 기다리고 있다.}%
}

%%K t08-apar-1 apar
\VUsaut{5.0mm}
\VUtitre{15.0pt}{60}{도표 제 8 호}
\VUsaut{3.0mm}
\VUcentre{13.2pt}{90}{\textit{첫째 마당.}}
\VUsaut{3.0mm}
\VUpk{10.2pt}{40}{거두기 (*).}
\VUsaut{4.0mm}
\VUpk{10.2pt}{40}{들의 모습.}
\VUsaut{3.0mm}

%%K t08-c1-01-1 p
1. --- \VUgras{여름}과 함께 들의 모습이 다시 한 번 바뀌었다. 고랑에 맡긴 씨가 싹텄다.
푸른 어린 싹이 땅에서 나와 이삭을 맺었다. 낟알이 여물고 익어, 이제는 거두기만 하면 된다.

%%K t08-c1-02-1 p
2. --- \VUgras{들꽃}이 피었다. \VUgras{수레국화} \textsuperscript{(1)},
\VUgras{개양귀비} \textsuperscript{(2)}, \VUgras{디기탈리스} \textsuperscript{(3)}가 밀밭의
한결같은 빛깔에 제 싱그러운 \VUgras{꽃부리}의 흥겨운 대비를 섞는다.

%%K t08-c1-03-1 p
3. --- 과일나무는 잎으로 덮였고 열매가 익었다. 작은 \VUgras{버찌나무}
\textsuperscript{(4)}에는 고운 \VUgras{버찌} \textsuperscript{(5)}가 그득한데, 아이들이 그것을
큰 즐거움으로 따 먹는다.

%%K t08-c1-04-1 p
4. --- 이제부터 가축은 온종일 바깥 공기 속에서 지낼 수 있다.

%%K t08-c1-04-2 p
\VUgras{새끼 양} \textsuperscript{(6)}들이 자라서 촘촘한 털가죽을 두른 고운 \VUgras{암양}
\textsuperscript{(7)}과 고운 \VUgras{숫양} \textsuperscript{(8)}이 되었다. 그것들은
\VUgras{목장} \textsuperscript{(9)}의 \VUgras{풀}을 조용히 뜯는다. 데이지로 알록달록해진
목장이다.

%%K t08-c1-04-3 p
곧 \VUgras{그} 짐승들의 털을 깎을 것이다. 그 \VUgras{양털}로 우리 옷의 나사를 짠다.

%%K t08-tit-2 sub
\VUsaut{5.0mm}
\VUpk{10.2pt}{40}{양치기.}
\VUsaut{3.4mm}

%%K t08-c1-05-1 p
5. --- \VUgras{양치기} \textsuperscript{(10)}는 그것들을 별로 눈여겨보지 않는 듯하다. 그는
지평선을 지나가는 뇌우만 걱정한다. 그리고 \VUgras{양 지키는 이} \textsuperscript{(13)}는
더욱 무심해 보인다 — 그는 \VUgras{호리병} \textsuperscript{(14)}을 입술로 가져가고 있다.

%%K t08-c1-06-1 p
6. --- 더위가 짓누른다. \VUgras{개} \textsuperscript{(15)}까지도 목을 축이러 오기 때문이다.
그것은 \VUgras{개울} \textsuperscript{(16)}의 시원한 물을 핥으며, 물가 풀 속에 숨어 있던
딱한 \VUgras{개구리} \textsuperscript{(17)}를 놀라게 한다. 개구리는 맑은 물로 뛰어드는데,
거기에는 \VUgras{수련} \textsuperscript{(18)}이 피어 있다.

%%K t08-c1-07-1 p
7. --- \VUgras{할미새} \textsuperscript{(19)} 한 마리가 \VUgras{샘} \textsuperscript{(20)} 곁에서
목욕한다. 두 바위 사이에서 솟는 샘이다. 그러고는 \VUgras{가시덤불} \textsuperscript{(21)}과
\VUgras{산사나무 덤불} \textsuperscript{(22)} 아래에 숨는다.

%%K t08-tit-3 sub
\VUsaut{5.0mm}
\VUpk{10.2pt}{40}{거두기.}
\VUsaut{3.4mm}

%%K t08-c1-08-1 p
8. --- 시골 아이들에게 밀 거두기는 아주 신나는 철이다. 그들은 즐겁게
\VUgras{뒹군다}\textsuperscript{(24)} — \VUgras{짚더미} \textsuperscript{(23)} 위에서. 그들은
\VUgras{거두는 일꾼} \textsuperscript{(26)}들을 돕는다. 일꾼들이 \VUgras{밀밭}
\textsuperscript{(27)}을 벨 때 — 제 \VUgras{큰 낫} \textsuperscript{(25)}으로,
\VUgras{숫돌} \textsuperscript{(30)}에 잘 갈아 둔 낫으로 — 아이들은 밀을 모아
\VUgras{밀단} \textsuperscript{(31)}이나 \VUgras{작은 단} \textsuperscript{(32)}으로 묶는다.

%%K t08-c1-09-1 p
9. --- 이따금 그들 가운데 가장 큰 아이에게 \VUgras{낫} \textsuperscript{(12)}으로
\VUgras{금빛 대} \textsuperscript{(28)}를 베게 허락한다. 그 대는 낟알이 그득한 묵직한
\VUgras{이삭} \textsuperscript{(29)}을 달고 있다. 작은 아이들은 \VUgras{가축}
\textsuperscript{(33)}을 지키면서 — 가축은 \VUgras{목장} \textsuperscript{(34)}에서
풀을 뜯는다, 큰 \VUgras{보리수} \textsuperscript{(35)} 그늘 아래에서 — 제 \VUgras{잠자리채}
\textsuperscript{(36)}로 고운 \VUgras{나비}를 잡는다.

%%K t08-c1-09-2 p
어떤 아이들은 \VUgras{수레} \textsuperscript{(37)}에 올라가, \VUgras{쇠스랑}
\textsuperscript{(38)}으로 건네주는 밀단을 가지런히 놓기도 한다.

%%K t08-c1-10-1 p
10. --- \VUgras{벌판} \textsuperscript{(39)} 전체에 활기가 넘친다 — 그 벌판에서
\VUgras{종달새} \textsuperscript{(41)}가 날아오른다. 모두가 거두기에 매달려 있다. 그러나
가난한 이들은 여전히 옛 연장을 쓴다 — \VUgras{큰 낫}, \VUgras{낫}, \VUgras{도리깨}
\textsuperscript{(40)}다. 넉넉한 지주들은 제 밀이나 제 \VUgras{호밀} \textsuperscript{(43)}을
\VUgras{베는 기계} \textsuperscript{(42)}로 베게 한다. 그러고는 밀단을 곡식 광으로 나르지
않고 그 자리에 큰 \VUgras{낟가리} \textsuperscript{(44)}를 쌓는다.

%%K t08-c1-11-1 p
11. --- 그때 \VUgras{타작 기계} \textsuperscript{(47)}를 끌어온다. 그것을
\VUgras{이동 증기기관} \textsuperscript{(45)}이 돌리는데, \VUgras{벨트} \textsuperscript{(46)}로
돌린다. 그렇게 하여 낟알이 \VUgras{짚}에서 갈라진다. 씨 뿌린 밭을 떠나지 않고서 말이다.

%%K t08-tit-4 sub
\VUsaut{5.0mm}
\VUpk{10.2pt}{40}{뇌우.}
\VUsaut{3.4mm}

%%K t08-c1-12-1 p
12. --- 거두는 철에는 자주 \VUgras{큰 뇌우} \textsuperscript{(53)}가 온다. 먼저 멀리서
\VUgras{천둥} 소리가 들린다. \VUgras{서풍의 돌풍} \textsuperscript{(54)}이 불기 시작하여
헐떡이며 \VUgras{숲} \textsuperscript{(49)}의 가장 굵은 나무들을 휘게 한다. 검은
\VUgras{먹구름} \textsuperscript{(56)}이 하늘을 덮치고, \VUgras{번개} \textsuperscript{(55)}의
눈부신 지그재그가 그것을 가른다. \VUgras{비}가, 때로는 \VUgras{우박}이 쏟아져, 벨 채비가
된 곡식을 쓰러뜨린다.

%%K t08-c1-13-1 p
13. --- 번개는 자주 건물에 불을 놓는다. 그렇게 지난해에는 \VUgras{풍차}
\textsuperscript{(50)}의 지붕이 잿더미가 되고 그 \VUgras{날개} \textsuperscript{(51)} 둘이
부서졌다.

%%K t08-c1-14-1 p
14. --- 몇 시간 뒤에 고요가 되살아난다. 지평선에 눈부신 \VUgras{무지개}
\textsuperscript{(52)}가 나타나고, 잠시 멈추었던 자연이 제 삶을 다시 시작한다.
\VUgras{오두막} \textsuperscript{(48)}으로 몸을 피했던 시골 사람들은 다시 일에 나서서,
방금 입은 피해를 씩씩하게 고치려 애쓴다.

%%K t08-tit-5 sub
\VUsaut{6.0mm}
\VUcentre{13.2pt}{90}{\textit{둘째 마당.}}
\VUsaut{3.6mm}

%%K t08-tit-6 sub
\VUpk{10.2pt}{40}{사냥. --- 포도 거두기.}
\VUsaut{1.4mm}
\VUpk{10.2pt}{40}{물고기 잡이.}
\VUsaut{4.0mm}

%%K t08-c2-01-1 p
1. --- 고장에 따라 \VUgras{팔월} 끝 무렵이나 \VUgras{구월} 중턱에 사냥이 열린다. 첫
햇살이 \VUgras{도마뱀} \textsuperscript{(2)}을 구멍에서 나오게 하자마자, \VUgras{사냥꾼}
\textsuperscript{(5)}이 길을 나선다 — 제 \VUgras{개} \textsuperscript{(4)}들과 함께.

%%K t08-c2-02-1 p
2. --- 그는 두꺼운 나사로 지은 단단한 \VUgras{사냥옷} \textsuperscript{(9)}을 입고, 가죽
\VUgras{각반}을 찼다. 그것으로 젖은 풀밭을 걸을 수 있는데, 거기에는 \VUgras{두꺼비}
\textsuperscript{(1)}가 숨어 있다. 그는 제 \VUgras{엽총} \textsuperscript{(6)}과
\VUgras{사냥 자루} \textsuperscript{(10)}를 들었다. 그리고 이제 떠난다.

%%K t08-c2-03-1 p
3. --- 쫓는 열기에 그는 포도밭 한가운데에 이르는데, 거기에서는 포도 거두기가 한창이다.
포도 농사꾼의 아이들은 여느 때 길에서 염소를 지키지만, 오늘은 아무 데나 풀을 뜯게
내버려 둔다. 그리고 늙은 \VUgras{숫염소} \textsuperscript{(3)}를 말뚝에 매어 둔 뒤 — 그
염소는 제 자유를 틈타 달아나기를 마다하지 않을 것이다 — 아이들도 저희 몫의 즐거움을
누린다. 하나는 \VUgras{따는 바구니} \textsuperscript{(12)}에서 포도 한 송이를 집어,
\VUgras{즙} \textsuperscript{(14)}을 \VUgras{잔} \textsuperscript{(13)}에 짜 넣어 새 즙(아직
익지 않은 포도주)을 만들며 논다. 다른 하나는 방금 엮은 \VUgras{잎 관}
\textsuperscript{(15)}으로 머리를 꾸민다.

%%K t08-c2-03-1 p suite
그들 곁에서는 뻔뻔한 \VUgras{참새} \textsuperscript{(16)}들이 포도알을 훔치러 압착기 앞까지
온다.

%%K t08-c2-04-1 p
4. --- 아이들이 노는 동안 어른들은 일한다. \VUgras{포도 따는 이} \textsuperscript{(18)}들은
\VUgras{등짐 바구니} \textsuperscript{(17)}로 포도를 술광까지 나른다. \VUgras{통메장이}
\textsuperscript{(19)}는 \VUgras{통} \textsuperscript{(20)}을 고치고, \VUgras{널}
\textsuperscript{(22)}과 \VUgras{밑판} \textsuperscript{(21)}이 성한지 살피며, 부러진
\VUgras{쇠테} \textsuperscript{(23)}를 갈아 끼운다. \VUgras{큰 통} \textsuperscript{(25)}도
그가 만들었다 — 그 안에 \VUgras{포도알} \textsuperscript{(26)}을 부어 익히는 것이다.

%%K t08-c2-05-1 p
5. --- 익기가 끝나면 포도주를 따라 내고, 남은 찌꺼기를 \VUgras{압착기}
\textsuperscript{(28)} 위에 올린다. \VUgras{큰 나사} \textsuperscript{(29)}로 조금씩 조이면서
즙을 짜내는데, 그 즙은 \VUgras{작은 통} \textsuperscript{(32)}으로 흐르고, 거기에서 다시
\VUgras{포도주} \textsuperscript{(31)}를 얻는다.

%%K t08-tit-8 sub
\VUsaut{6.0mm}
\VUpk{10.2pt}{40}{맥주와 사과술.}
\VUsaut{3.4mm}

%%K t08-c2-06-1 p
6. --- \VUgras{술광} \textsuperscript{(27)} 벽을 \VUgras{홉} \textsuperscript{(30)} 덩굴이
타고 오른다. 여기에서 그것은 꾸밈용 풀일 뿐이다. 그러나 포도가 아주 드문 몇몇 나라에서는
\VUgras{맥주}를 만들려고 홉을 가꾼다.

%%K t08-c2-07-1 p
7. --- 작은 \VUgras{사과나무} \textsuperscript{(33)}에는 사과가 그득한데, 익으면 훌륭한
\VUgras{사과술}을 낼 것이다. 제 \VUgras{새장} \textsuperscript{(34)} 안에서
\VUgras{카나리아} \textsuperscript{(35)}와 \VUgras{방울새} \textsuperscript{(36)}가 마음껏
뛰노는 참새들을 부러운 눈으로 바라본다.

%%K t08-c2-08-1 p
8. --- 눈이 닿는 데까지, 언덕 비탈에 \VUgras{포도 받침대} \textsuperscript{(40)}가 줄지어
서서 훌륭한 \VUgras{포도 그루} \textsuperscript{(39)}를 받치고 있다. 그 그루에는 묵직한
\VUgras{송이}가 달려 있다.

%%K t08-c2-08-2 p
한 해 내내 \VUgras{포도 농사꾼} \textsuperscript{(42)}은 제 \VUgras{포도밭}
\textsuperscript{(38)}을 돌보느라 일했고, 이제 고운 수확으로 그 수고를 보상받는다.
\VUgras{포도 따는 아낙} \textsuperscript{(41)}들이 수레에 올린 통을 채우고 —
\VUgras{노새} \textsuperscript{(43)}가 끄는 수레다 —, 모두가 즐겁다.

%%K t08-tit-9 sub
\VUsaut{5.0mm}
\VUpk{10.2pt}{40}{물고기 잡이.}
\VUsaut{3.4mm}

%%K t08-c2-09-1 p
9. --- \VUgras{낚시꾼} \textsuperscript{(44)}들도 마찬가지다. 그들은 아침부터 제
\VUgras{낚싯대} \textsuperscript{(45)}를 들고 물가에 자리를 잡았다.

%%K t08-c2-09-2 p
\VUgras{물고기} \textsuperscript{(47)}는 \VUgras{낚싯바늘}을 문다 — 그들이 제
\VUgras{낚싯줄} \textsuperscript{(46)}을 물에 담그자마자. 그리고 \VUgras{미끼}를 다시 끼울
겨를도 없이 새 먹이가 줄 끝에서 파닥거린다.

%%K t08-c2-09-3 p
그래도 \VUgras{그물 던지는 이} \textsuperscript{(48)}가 더 많이 잡는다. 그는 제
\VUgras{배} \textsuperscript{(49)}에서 \VUgras{던짐그물}을 \VUgras{강} \textsuperscript{(50)}
한가운데로 던진다.

%%K t08-c2-10-1 p
10. --- 가을은 좋은 나들이의 철이다. 걸어서, \VUgras{말 타고} \textsuperscript{(52)},
\VUgras{자전거로} \textsuperscript{(53)}, \VUgras{자동차로} \textsuperscript{(54)}.
\VUgras{길} \textsuperscript{(51)}이 깨끗하고, 사람들은 마지막 좋은 날들을 누린다. 벌써
\VUgras{제비} \textsuperscript{(56)}들이 지붕 위에 모여 온 날개로 더 따뜻한 나라를 향해
날아오르려 하기 때문이다. 늙은 \VUgras{호두나무} \textsuperscript{(57)}의 잎이 누레지고
\VUgras{호두}가 떨어진다. 키 큰 \VUgras{나무}들도 곧 잎을 떨굴 것이다 —
\VUgras{꽃다발} \textsuperscript{(58)}처럼 모여 선 나무들, \VUgras{언덕}
\textsuperscript{(59)} 위의 나무들이다.

%%K t08-c2-11-1 p
11. --- \VUgras{해} \textsuperscript{(60)}가 늦게 뜨고 낮이 줄어들며, \VUgras{아침}에는
제법 살을 에는 추위가 느껴지기 시작한다. 그리고 벌써 \VUgras{멧도요}
\textsuperscript{(61)} 떼가 우리의 손님 대접 없는 기후를 떠나고 있다.

\VUsaut{6.0mm}
\VUfilet{22mm}
